% Source: https://artofproblemsolving.com/community/c3h1672490_system_of_equations
% Stable identity: aops-topic-1672490-q01
\textbf{Combine into one complex equation.} Let $z=b+ci$. Then
\[
z^3=b^3+3b^2(ci)+3b(ci)^2+(ci)^3=\left(b^3-3bc^2\right)+\left(3b^2c-c^3\right)i ,
\]
so for real $b,c$ the system says exactly $z^3=2+2i$.

\textbf{Solve $z^3=2+2i$.} In polar form $2+2i=2\sqrt2\left(\cos45^\circ+i\sin45^\circ\right)$. Its cube roots have modulus $\left(2\sqrt2\right)^{\dfrac{1}{3}}=\sqrt2$ and arguments $15^\circ$, $15^\circ+120^\circ=135^\circ$, $15^\circ+240^\circ=255^\circ$:
\[
z=\sqrt2\left(\cos\theta+i\sin\theta\right),\qquad\theta\in\{15^\circ,135^\circ,255^\circ\}.
\]
Using $\cos15^\circ=\dfrac{\sqrt6+\sqrt2}{4}$ and $\sin15^\circ=\dfrac{\sqrt6-\sqrt2}{4}$, and $\cos255^\circ=-\sin15^\circ$, $\sin255^\circ=-\cos15^\circ$:
\[
\boxed{(b,c)=\left(\dfrac{\sqrt3+1}{2},\dfrac{\sqrt3-1}{2}\right),\quad(-1,1),\quad\left(\dfrac{1-\sqrt3}{2},\ -\dfrac{1+\sqrt3}{2}\right).}
\]
Check $(-1,1)$: $-1+3=2$ and $3-1=2$. Check the first: $b^2=\dfrac{2+\sqrt3}{2}$, $c^2=\dfrac{2-\sqrt3}{2}$, $b^2-3c^2=-2+2\sqrt3$, so $b(b^2-3c^2)=\dfrac{(\sqrt3+1)(2\sqrt3-2)}{2}=2$; similarly $c(3b^2-c^2)=\dfrac{(\sqrt3-1)(2+2\sqrt3)}{2}=2$. Since the map $(b,c)\mapsto z$ is one-to-one and $z^3=2+2i$ has exactly three roots, these are all the real solutions.
